\VerdierPhysicalPageAnchor{0023}%
$\mathcal A$ possède suffisamment d'objets injectifs. Le foncteur dérivé
total permet donc, en particulier, de reconstruire les ``hyperdérivés à droite
de $F$'' ; mais il fournit plus d'informations. Par exemple, lorsqu'on a un
foncteur additif $G:\mathcal B\to\mathcal C$ entre catégories abéliennes, on
peut composer les foncteurs dérivés totaux :

\VerdierSemanticAnchor{intro-3-derived-composition}
\[
\mathrm R G\mathbin{\cdot}\mathrm R F:
\mathbf D^*(\mathcal A)\longrightarrow\mathbf D(\mathcal C).
\]
Dans les cas habituels, le foncteur
$\mathrm R G\mathbin{\cdot}\mathrm R F$ est canoniquement isomorphe au
foncteur dérivé total de $GF$ :

\VerdierSemanticAnchor{intro-3-derived-composition-isomorphism}
\[
\mathrm R(GF)\xrightarrow{\sim}
\mathrm R G\mathbin{\cdot}\mathrm R F,
\]
formule de composition qui traduit dans notre langage
\emph{la suite spectrale des foncteurs composés}, et qui est d'ailleurs plus
précise, car elle ne se borne pas à donner certaines informations sur les
objets de cohomologie de $\mathrm R GF(X)$
($X$ objet de $\mathbf D^*(\mathcal A)$), mais elle décrit entièrement le
complexe à ``quasi-isomorphisme près'' qui donne naissance à ces objets de
cohomologie ; complexe dont la connaissance permet, par exemple, de composer
à nouveau avec un autre foncteur dérivé total.

\VerdierSemanticAnchor{intro-4}
\noindent\textbf{4.}\quad
D'une manière générale, on est amené, dans ce formalisme, à travailler
directement avec les complexes à ``quasi-isomorphisme près'', complexes qu'on
se garde bien de remplacer par les seuls objets de cohomologie. (Si, le plus
souvent, les ``invariants'' que l'on recherche, au terme du calcul ou du
raisonnement, sont les objets de cohomologie des complexes, il n'en est pas
toujours ainsi : les différentielles des suites spectrales sont, elles aussi,
des ``invariants'' intéressants).

\VerdierSemanticAnchor{intro-4-derived-relations}
On établit, au niveau des complexes à quasi-isomorphisme près, ou au niveau des
foncteurs dérivés totaux, ou plus généralement au niveau des foncteurs entre
les catégories dérivées, diverses relations (relation d'isomorphie, relation
d'adjonction pour les théorèmes du type dualité, etc.). Ces relations
fourniront, par des procédés purement automatiques (les suites spectrales), des
informations sur les objets de cohomologie des complexes, sur les filtrations
naturelles dont ces objets sont munis, etc.

\VerdierSemanticAnchor{intro-4-technical-intermediaries}
Même lorsqu'on ne s'intéresse qu'aux valeurs des foncteurs dérivés
traditionnels, \emph{i.e.} aux seuls objets de cohomologie, les relations dans
les catégories dérivées sont souvent des intermédiaires techniques
indispensables pour arriver à étudier ces foncteurs et formuler pour eux
certaines relations importantes.

\VerdierSemanticAnchor{intro-4-simpler-relations}
Enfin, avantage essentiel, les relations qu'on établit dans les catégories
dérivées sont d'apparence et de fait plus simples que les relations qu'on en
